% Generated draft: requires mathematical and source review.
\documentclass[11pt]{article}
\usepackage[margin=25mm]{geometry}
\usepackage{fontspec}
\setmainfont{Latin Modern Roman}
\usepackage{amsmath,amssymb,mathtools,unicode-math}
\setmathfont{Latin Modern Math}
\usepackage{enumitem,needspace,xurl,hyperref}
\setlength{\parindent}{0pt}
\setlength{\parskip}{5pt}
\setlength{\emergencystretch}{3em}
\newcommand{\N}{\mathbb{N}}
\newcommand{\Z}{\mathbb{Z}}
\newcommand{\Q}{\mathbb{Q}}
\newcommand{\R}{\mathbb{R}}
\newcommand{\C}{\mathbb{C}}
\newcommand{\F}{\mathbb{F}}
\newcommand{\A}{\mathbb{A}}
\newcommand{\PP}{\mathbb{P}}
\newcommand{\E}{\mathbb{E}}
\newcommand{\eps}{\varepsilon}
\newcommand{\dd}{\,\mathrm{d}}
\begin{document}
\section*{UnsolvedMath ID: 11000316}
\textbf{Source problem code:} catalogue entry
\subsection{Definitions and mathematical statement}
Let \(F_n\) be the rank-\(n\) free group and \(V\) a specified \(\Q\)-vector space with a linear \(\operatorname{GL}(n,\Q)\)-action. Abelianization gives maps
\[
 \operatorname{Aut}F_n\longrightarrow\operatorname{GL}(n,\Z)
 \subset\operatorname{GL}(n,\Q),\qquad
 \operatorname{Out}F_n\longrightarrow\operatorname{GL}(n,\Q),
\]
through which \(V\) becomes a coefficient module for either group. The problem is to compute
\[
 H^k(\operatorname{Aut}F_n;V),\qquad
 H^k(\operatorname{Out}F_n;V)
\]
for the intended varieties of such modules. The source does not fix a module family, finite-dimensionality or algebraic-representation requirement, degree range, rank range or stable versus unstable scope. Nontrivial module action is part of the target, not replaced by trivial-coefficient cohomology.

It specifically asks how Looijenga's mapping-class-group result [62] might generalize to these contexts and cites Kawazumi [49] and Satoh [97] for recent progress. The exact Looijenga theorem and those results are absent, so no analogous formula or vanishing range is invented. A following well-known problem is announced but not supplied.
\subsection{Short English statement}
Compute twisted cohomology of free-group automorphism and outer automorphism groups for specified \(\operatorname{GL}(n,\Q)\)-modules, and seek the intended analogue of Looijenga's mapping-class result.
\subsection{Sources}
\begin{itemize}
\item S. Morita, chapter in \emph{Problems on Mapping Class Groups and Related Topics}, edited by Farb, Problem 4.10, supplied PDF page 366. \url{https://www.math.uchicago.edu/~farb/papers/mcgbook.pdf}. Recorded access 2026-07-29; extraction pdf-font-structure.
\item Looijenga [62], Kawazumi [49] and Satoh [97], with the roles described in the source; full bibliography and theorem statements are absent. The supplied summary loosely mentions Galatius and Johnson-type filtration work without an exact reference.
\item Local UnsolvedMath registry, \texttt{data/unsolvedmath/problems.json}, permanent ID \(11000316\), exact source code \texttt{catalogue entry}; supplied OPEN-TRIAGE summary reports partial related work.
\end{itemize}
\end{document}
